% !TEX root = ../main_zh.tex
\chapter{劳动力市场与失业}
\label{ch:labor}

在宏观经济学家关注的所有数字里，失业率大概是离研讨室最远的一个。GDP 增长率多一个百分点是一种抽象，而一个找不到工作的人却是实实在在的。在许多民主国家，失业率被读作对政府的一纸判决，从利率决策到延长失业救济，一长串政策都或明或暗地与它挂钩。本章要搭建的，正是诚实地读懂这个数字所需要的一整套工具：劳动年龄人口如何被分门别类，谁算失业、谁不算；为什么离开劳动参与率单看失业率几乎毫无意义；以及为什么一个健康的经济体应当存在正的失业，而不是把失业降到零。

思路是直截了当的。失业率是一个比率，我们必须先说清楚分子和分母里到底装着什么，所以第一步是对人口做划分。随后我们定义两个核心比率——失业率与劳动参与率——并说明为什么二者必须放在一起读。会计口径理清之后，我们按成因把失业分成三类：摩擦性失业、结构性失业与周期性失业。这一分类引出本章的核心概念——自然失业率（natural rate of unemployment），它进而界定了经济学家所说的“充分就业”以及经济体的潜在产出（potential output）。全程中我们都会留意那些制度与数据质量上的问题，正是它们让这些数字在实践中远比理论上更难解读，在中国的语境下尤其如此。

\section{人口的划分}

要把“失业”变成一个可度量的量，我们必须先确定谁才有资格被算作失业者。官方统计分两步建立这一划分，依次问两个问题：这个人\emph{能不能}工作？如果能，那他\emph{愿不愿意}工作？

第一层划分是看\textbf{有无劳动能力}。人口中有一部分人之所以在劳动力市场之外，不是出于选择，而是出于处境或规则：儿童、被收容者、全日制学生、现役军人、在押人员。无论他们自己愿不愿意，这些群体在第一层就被排除在外，因为他们并不能自由地去市场上谋一份工作。

在有劳动能力的人当中，第二层划分是看\textbf{有无劳动意愿}。一个既有能力又愿意工作的人，被算入\emph{劳动力}；一个有能力工作、却不愿意工作的人——选择不去领薪工作的家庭主妇、靠积蓄提早退休的人——则被算作\emph{非劳动力}。因此，劳动力是两个条件的交集：既有能力，又有意愿。

\defn{劳动力及其构成}{
在劳动年龄、未被收容的人口中：
\begin{itemize}
    \item \emph{劳动力}是所有既\emph{有能力}工作、又\emph{有意愿}工作的人。它进一步分为\emph{就业者}与\emph{失业者}。
    \item \emph{失业者}是那些有能力、也有意愿工作——正在积极寻找工作——但当下没有工作的人。
    \item \emph{非劳动力}指有能力工作、但在现行条件下不愿意工作的人。这一群体包括家庭主妇等创造了价值却没有产生可计量市场产出的人，也包括\emph{灰心丧气的劳动者}（discouraged workers）。
\end{itemize}
}

这一划分有两个特征值得强调，因为官方口径恰恰会在这两处误导我们。

第一，失业者并不等于“没有工作的人”。退休老人和幼儿同样没有工作，可两者都不算失业：退休老人是不愿意，幼儿是没能力。要被算作失业，一个人必须同时越过两道门槛——有能力且有意愿——并且仍然没有工作。失业是想工作的意愿与得到工作的机会之间的落差。

第二，灰心丧气的劳动者卡在边界上，恰好暴露了这套口径的一处微妙弱点。所谓灰心丧气的劳动者，是指那些原本愿意接受一份工作、却因确信无工可寻而放弃了主动寻找的人。按照调查实际采用的意愿标准——它以“是否在积极找工作”为关键——灰心丧气的劳动者被划归非劳动力，而不是失业者。于是他们同时从失业率的分子和分母里消失了。当经济下行把大量求职者推入灰心丧气的状态时，所度量的失业率反而可能下降，尽管劳动力市场实际上恶化了——原因正是这些人被重新归类到了劳动力之外。这也是为什么失业率不能孤立地读，我们下文还会回到这一点。

\rmkb[计量失业的两种口径，以及中国为何切换]{
中国长期使用城镇\emph{登记}失业率，它只统计那些走进当地劳动部门、为领取救济而登记为失业的人。由于登记本身既不完整、又受激励驱动，这一口径对真实的劳动力市场松弛程度反映得很差。中国后来开始按国际通行做法、依据住户调查逐月公布城镇\emph{调查}失业率，从而同时提升了数据的可靠性与国际可比性。一个更宽泛的提醒值得直白说出：中国的劳动力市场数据历来弱于相应的产出和价格数据，官方规划中的就业目标往往定得偏稳健（即偏低），以确保能够可靠地完成。在就业数据内部，两个群体——工人与应届大学毕业生——在政治上最为敏感，也最受密切关注。
}

\rmkb[失业与政治周期]{
不同经济体赋予失业的分量差别很大。在许多西方民主国家，失业率是核心的选举与政策变量，大量宏观经济政策都与它挂钩。中国的政策框架对失业数字本身的反应要间接得多，更多依靠其他大的实体经济指标；但这并不意味着就业在那里不重要，因为劳动力市场状况与其他每一个宏观经济问题都紧密交织。差别在于决策者盯着哪一个刻度盘，而不在于就业重不重要。
}

\section{两个核心比率}

划分一旦固定，用来概括它的两个比率就只是简单的比值了。第一个度量劳动力\emph{内部}的松弛程度，第二个度量劳动力相对于原则上能供给劳动的人口有多大。

\defn{失业率}{
\emph{失业率}是劳动力中处于失业状态的比例，
\[
u = \frac{\text{失业人口}}{\text{劳动力人口}}
  = \frac{\text{失业人口}}{\text{失业人口} + \text{就业人口}}.
\]
它的分母是劳动力，\emph{而非}全部人口：不愿意或没能力工作的人在分子和分母里都被排除在外。
}

\defn{劳动参与率}{
\emph{劳动参与率}（labor-force participation rate）是有劳动能力的人口中实际进入劳动力的比例，
\[
\mathrm{LFPR} = \frac{\text{劳动力人口}}{\text{劳动年龄人口}}.
\]
它度量的是，在有能力参与劳动的人当中，有多少比例\emph{愿意}去劳动。
}

这两个比率回答的是不同的问题，而且各自都对另一个所看见的东西视而不见。失业率只往劳动力内部看，它对“有多少有能力的人选择待在劳动力之外”只字不提；劳动参与率只数有多少人进了劳动力，它对“进来的人过得怎样”只字不提。任何一个单独报出来，都可以被读成截然相反的意思。

\state{看失业率一定要参考劳动参与率}{
失业率下降只有在劳动参与率稳住的前提下才是好消息。如果人们离开劳动力——提前退休、重返校园，或干脆放弃寻找而成为灰心丧气的劳动者——那么 $u$ 的分子和分母会同时缩小，于是即便劳动力市场正在恶化，失业率也可能下降。\emph{失业率低并不意味着劳动力供给充足。}
}

最清楚的例证是一个福利和转移支付都很优厚的经济体。当不工作也过得舒服时，许多有能力的劳动者会选择不参与；劳动参与率下降，而由于灰心丧气者和自愿赋闲者都不在劳动力之内，所度量的失业率也随之下降。然而结果不是劳动力市场的繁荣，而是它的萎缩——美国就曾呈现出恰好这样的格局：劳动参与率下行，失业率却很低，同时又抱怨劳工不足。失业率看上去很强劲，劳动参与率却揭示出，可供使用的劳动力池子已经变薄了。

\subsection{非农就业}

基于调查的比率并非唯一会撼动市场的劳动力市场数据。在美国，最受密切关注的单项发布是\emph{非农就业}（nonfarm payrolls，NFP），即农业部门之外有薪工作岗位数量的月度变化。

\defn{非农就业（NFP）}{
\emph{非农就业}统计的是经济中除农业部门以外的受薪雇员数量。其头条数字是环比月度变化，被解读为净新增（或净流失）的就业岗位。
}

农业被刻意剔除。农业就业具有强烈的季节性和波动性，随播种与收获的日历起伏，这些起伏对劳动力市场的根本健康状况几乎说明不了什么。把它去掉之后，留下的是对“经济到底创造或损失了多少岗位”更干净的读数。由于非农就业是对岗位创造既及时又直接的度量，一次大幅偏离预期的发布，可能给金融市场带来巨大的冲击。

\section{为什么劳动力市场对“人”可能无法出清}

我们现在从度量转向机制：\emph{为什么}会存在失业？在一个理想化的竞争性市场里，过剩会把价格压低，直到供给量等于需求量，没有任何东西卖不出去。失业的持续存在告诉我们，劳动力市场并不会如此干净地出清。作为一条组织原则：\emph{劳动力市场上出现失业，说明市场上存在摩擦（friction）。}三种失业，就是三种不同类型的摩擦。

\subsection{摩擦性失业}

第一种摩擦很简单，就是把一个劳动者匹配到一份工作需要时间。\emph{摩擦性失业}（frictional unemployment）是人们在寻找并被分配到岗位的过程中出现的失业，即便在一个其他方面都健康的经济体里也存在。

\defn{摩擦性失业}{
\emph{摩擦性失业}是因匹配劳动者与工作岗位需要时间而产生的失业。劳动者出于与自身能力无关的原因失去或离开工作，而要找到合适的新匹配——合适的企业、职业或地点——需要时间去寻找和谈判。
}

有些摩擦性失业是个人层面的：处于两份工作之间的劳动者正经历一段匹配期，不断比较各种录用机会，直到合意的那个出现。有些则是行业层面的：当某个行业重组、技术变迁，或经济活动从一个地区迁往另一个地区时，劳动者必须跨企业、跨职业或跨地点流动，而这种再配置无法瞬间完成。在这个意义上，摩擦性失业是一个不断创造与毁灭岗位的动态经济不可避免的副产品。它不是需求疲软的症状——劳动者失去工作的原因与他们是否胜任毫无关系——而是\emph{寻找与匹配需要时间}这一无法消除的事实的症状。

\subsection{结构性失业}

劳动力市场是一个特殊的市场，因为它交易的商品依附于一个人，而人并不会像一蒲式耳小麦那样被压价到使市场出清的工资水平。\emph{结构性失业}（structural unemployment）是因为工资被维持在高于市场出清水平之上而持续存在的失业：此时劳动供给量超过需求量，一部分愿意工作的人找不到工作。它有三个经典来源。

\defn{结构性失业}{
\emph{结构性失业}是因工资被维持在其市场出清水平之上、从而造成劳动长期供过于求而持续存在的失业。它的三个标准来源是工会、最低工资立法与效率工资。
}

\paragraph{工会与集体谈判。}
工会代表其成员进行集体谈判，谈判桌上的头号议题就是涨工资。从一个本来会在均衡工资处出清的市场出发，一个被谈判（或被行政力量强加）抬到该水平之上的工资，提高了那些保住饭碗者的报酬，却减少了企业愿意提供的岗位数量，从而让一部分人陷入失业。实际上，工会保护的是它的\emph{内部人}——那些在更高工资下仍然就业的成员——代价则由那些被高工资挤出工作的外部人承担。

\paragraph{最低工资。}
立法规定的最低工资是同一机制的政府版本。如果它被设定在低技能劳动的市场出清工资之上，就会在提高保住饭碗者收入的同时，把生产率最低的劳动者挤出就业——一项出发点良好、却可能办坏事的政策。这正是民粹主义的陷阱：一项听上去纯粹是体恤劳动者的措施，由于忽视了企业会如何反应，反而可能让最脆弱的劳动者处境更糟。它提出了一个真正棘手的问题——多高的工资下限或福利支持才算合理？——对此，市场力量给出一个答案，政治压力给出另一个答案。

\paragraph{效率工资。}
第三个来源更为微妙，因为在这里是企业\emph{主动}抬高工资。按照\emph{效率工资}（efficiency wage）理论，一家对工人努力程度掌握信息不完备的企业，可能会发现支付高于市场出清水平的工资反而有利可图，因为这能提高生产率：工资溢价降低了人员流动、吸引了更优秀的应聘者，并让现有工人有了“怕失去”的东西，于是他们更卖力、更少偷懒。企业以更高的工资支出换取每个工人更高的努力。但同样这份买来努力的溢价，也意味着工资位于出清水平之上，因此从总量上看，被雇佣的工人少于竞争性工资下本应雇佣的数量——而这部分过剩，便以失业的形式显现出来。

\rmkb[结构性失业关乎“人”，而非商品]{
三个来源——工会、最低工资与效率工资——共有一个根源：劳动力市场交易的是活生生的人，而不是匿名的货物。工资向下具有黏性，合同与法律对其构成约束，公平感与士气会影响努力程度，而关于谁在卖力工作的信息又是不完备的。这些“人”的特征，恰恰是一个无摩擦的商品市场所缺少的，也正是劳动力市场不会简单出清的原因。
}

\subsection{周期性失业}

前两类失业是经济正常运转的固有属性，第三类则是经济周期的属性。\emph{周期性失业}（cyclical unemployment）是随宏观经济总体状况起落而升降的失业：衰退收缩了对劳动的需求、把工人推出岗位，繁荣则反其道而行之。

\defn{周期性失业}{
\emph{周期性失业}是由经济周期中总需求波动所引起的那部分失业。它在衰退中为正——此时对劳动的需求下降；在经济过热的繁荣中则可能为负。
}

在繁荣时期，周期性失业实际上可以转为负值：企业招不到足够的工人，雇员可能被迫加班，工作量超过他们自愿选择的水平。不过这种过热通常被认为是暂时的，而非经济的一种持久状态。

\section{自然失业率、充分就业与潜在产出}

现在我们可以把三类失业组装成本章的核心组织思想了。摩擦性失业与结构性失业是\emph{长期}特征：它们反映了劳动力市场永久性的构造——寻找摩擦、制度、信息——即便在一个完全没有需求冲击的经济体里也会持续存在。相比之下，周期性失业是\emph{短期}成分，是随繁荣与衰退来去的那一部分。自然失业率把长期的那一部分单独剥离出来。

\defn{自然失业率}{
\emph{自然失业率}是摩擦性失业与结构性失业之和：
\[
u^{*} = u_{\text{摩擦}} + u_{\text{结构}}.
\]
它是周期性失业为零时所通行的失业率。由于它由长期的、结构性的因素所驱动，自然失业率无法被短期宏观经济政策所撼动。
}

自然失业率就是把经济周期“关掉”之后剩下来的东西。一个有用的看法是：当\emph{实际} GDP 等于\emph{潜在} GDP 时，外生需求冲击基本为零，周期性失业为零，于是所观测到的失业率恰好就是自然失业率——摩擦性失业与结构性失业之和。既然自然失业率由长期力量所决定，短期稳定政策就无法把它降下来；短期政策真正应当瞄准的，只有\emph{周期性}那一部分。

这就使经济学家所说的“充分就业”这个容易被误读的词变得更加清晰。

\state{充分就业指周期性失业为零，而非失业率为零}{
\emph{充分就业}是周期性失业为零、所度量的失业率等于自然失业率 $u^{*}$ 的状态。它\emph{并不}意味着失业率为零。即便处于充分就业，也仍然存在摩擦性失业（人们仍在两份工作之间流动）和结构性失业（市场的某些部分工资仍高于出清水平）。一个所度量失业率为零的经济体并不健康，那将是一个劳动者完全没有余地在岗位之间流动的经济体。
}

充分就业在产出上的对应物是\emph{潜在 GDP}。

\defn{潜在产出与潜在增长率}{
\emph{潜在 GDP} 是经济体在充分就业时所生产的产出水平——也就是当实际产出剥离了短期冲击、失业处于其自然率时的产出。\emph{潜在增长率}是去除短期波动之后经济体的长期增长率，由资本积累和技术进步所驱动。
}

\rmkb[潜在产出与其说是测量，不如说是一种信念]{
潜在产出与潜在增长率都是理论建构：我们从未观测到一个被干净地剥去了冲击的经济体。在实践中，对潜在增长率的估计更接近于一种\emph{信念}，通常是通过总结过去的经验、再向前外推而形成的。当一个政府宣布——比如说——要在某一年之前把经济总量翻一番（折算下来约为每年平均百分之五的增速）时，它表达的是关于潜在增长率的一种看法，而不是在报告一个测量到的事实。自然失业率也具有同样的地位：它真实而重要，但它是被推断出来的，而不是从某个刻度盘上读出来的。
}

把时间维度归纳一下：摩擦性失业与结构性失业是长期成分，由劳动力市场的结构所固定，短期政策无从触碰；周期性失业是短期成分，是需求管理政策唯一有望撼动的那一部分。充分就业则是那个短期成分已被压到零、只剩自然失业率的那一点。

\section{人口结构与劳动参与：两段旁论}

劳动参与率引出的问题会自然地通向人口学，从中国的经验里有两点观察值得记下来，尽管它们都不能干净利落地纳入上文的失业会计框架。

\rmkb[为什么中国女性的劳动参与率如此之高]{
尽管有着女性相夫教子的悠久文化传统，中国的女性劳动参与率却位居世界前列。有两股力量朝同一个方向推动。第一是\emph{家庭收入}水平：在单一劳动者无法养活一家人的地方，夫妻双方都必须工作才能勉强维生，参与与否其实算不上一种选择。第二是\emph{技术变迁}：随着生产从体力劳动转向脑力劳动，女性所能从事的工作种类、以及她们从事这些工作的生产率，都与男性趋于接近，旧有的职业分工因而被侵蚀。当体力劳动越来越不重要、脑力劳动越来越重要时，男性体力上的历史溢价便逐渐消退。
}

同样这种向脑力劳动的转移也重塑着生育，而这里的经济学值得写下来。考虑一个在消费 $C$ 与孩子数量 $N$ 之间做选择的家庭，其效用为 $U(C, N)$，预算约束为
\[
C + N P = w\,(1 - t N),
\]
其中 $w$ 是工资，$P$ 是抚养一个孩子的物品成本，$t$ 是抚养每个孩子所需的时间。$w\,(1 - t N)$ 一项是家庭在扣除了为抚养 $N$ 个孩子而从工作中抽走的时间之后的收入：每个孩子占用父母工作时间的比例为 $t$，因此每个孩子除了直接成本 $P$ 之外，还附带 $w t$ 的\emph{机会成本}，即放弃的工资。

\rmkb[数量与质量的权衡，以及人口转型]{
随着工资 $w$ 上升来读这条预算约束，便勾勒出了人口转型的轨迹。在工业化的早期阶段，$w$ 迅速增长，而孩子的时间成本仍然不高，于是家庭对收入提高的反应是生\emph{更多}的孩子。当 $w$ 进一步上升，每个孩子的时间成本 $w t$——即父母放弃的工资——开始占据主导，孩子在“随收入而水涨船高”的那一种资源上变得昂贵起来：父母的时间。于是家庭转向少生孩子、在每个孩子的\emph{人力资本}（human capital）上投入更多；这正是经典的\emph{数量与质量的权衡}（quantity--quality tradeoff）。随之而来的生育率下降便是人口转型，它意味着人口老龄化不是政策的意外，而是一个建立在脑力劳动之上的社会在发展过程中基本无法回避的一个阶段。
}
